% Source: ../蒋旭个人简历202606(1).pdf
% Personal details, dates, experiment settings and results follow the source PDF.

\section{教育背景}
\begingroup
\fontsize{10}{14}\selectfont
\noindent\begin{tabularx}{\linewidth}{@{}X r@{}}
  \textbf{北京大学} -- 软件与微电子学院 -- 电子信息 -- 硕士研究生（保研） &
  {\color{SecondaryText}2025-09 $\sim$ 2028-06}\\[2pt]
  \textbf{北京大学} -- 信息科学技术学院 -- 应用物理 -- 本科 &
  {\color{SecondaryText}2021-09 $\sim$ 2025-06}\\[1pt]
  \hspace*{4.8em}国家发展研究院 -- 经济学双学位 &
\end{tabularx}\par
\endgroup
\vspace{3pt}
\noindent\begin{tabularx}{\linewidth}{@{}l @{\hspace{0.65em}} X@{}}
  \textbf{主修课程} &
  算法分析与设计；自然语言处理；数据结构与算法；大语言模型后训练实践；数字图像处理
\end{tabularx}

% !TeX root = resume.tex
% Per the user's request, this draft copies only the JD internship section from
% ../陈凯权_简历.pdf. The role, dates, responsibilities and metrics below follow
% that reference and can be personalized in this file.

\section{实习经历}
\companyheading{assets/logos/jd-joy-rounded.pdf}{京东集团\,·\,京东科技}{算法开发}{2026.06 -- 至今}
\projectheading{工具调用大模型训练与推理优化}{}
\begin{itemize}[label=\textbullet,leftmargin=1.2em,itemsep=3pt,topsep=0pt]
  \item \textbf{81B 工具调用模型全参微调：}基于定制 LLaMA-Factory，对内部 Llama 系稠密模型开展多机多卡全参数 SFT；采用 BF16、DeepSpeed ZeRO-3 参数／优化器 CPU Offload、FlashAttention 2 与梯度检查点，支持 16K 序列训练；通过 Neat Packing 与块对角注意力隔离打包样本，适配多工具调用、工具返回及最终回答模板。
  \item \textbf{数据配比与 Chat／工具路由消融：}对比通用 FC、业务及长多轮混合数据，评估工具调用与无需调用工具的边界。通用 FC 训练使 Non-Live／Live 准确率分别由 68.77\%／50.56\% 提升至 \textbf{90.12\%／82.68\%}，同时识别不调用工具能力退化；通用与业务数据联合训练的业务评测平均准确率为 \textbf{93.23\%}，较仅业务数据的 91.13\% 提升 \textbf{2.10 个百分点}。
  \item \textbf{MTP 后训练与推测解码：}冻结已微调的 81B 主模型，后挂单层 Llama Decoder 与特征融合层，复用冻结的词嵌入及输出头，训练轻量单步 Drafter，无需重训主模型。家居、万能博士、外卖场景下，MTP 单独相对原模型加速 \textbf{1.26--1.49 倍}，与 FP8 组合加速 \textbf{1.65--1.94 倍}；对比跨场景接受率，评估业务特化收益与泛化边界。
  \item \textbf{智能家居工具调用数据体系：}针对中文数据稀缺与 Teacher 标签幻觉，对比开源数据及近 99.2 万条业务日志，制定“HA 运行时真值为主、开源增强、业务校准”的路线；以 25 套家庭拓扑、实体状态与真实执行结果构造标签，按设备能力暴露工具，将 Teacher 限定为生成用户表达，解耦语言多样性与调用标签正确性。
  \item \textbf{数据规模化与可执行验证：}以确定性真值批量生成多种表达，结合 Qwen／GLM 双 Teacher、Schema 约束、语义去重及 64 路并发与失败恢复；完成多设备控制、查询及拒绝场景的中文增强和协议归一，通过静态规则、25 套 HA 环境运行时测试、分阶段模型审计与人工抽样，形成\textbf{约 9.8 万条}可用于 SFT 的数据（79,733 条核心数据、17,856 条增强数据）。
\end{itemize}


\section{项目经历}
\projectheading{基于 PyTorch 实现的轻量级大语言模型复现与训练}{2026-03 $\sim$ 2026-04}
\textbf{项目介绍：}为深入研究 LLM 内部原理和机制，纯 PyTorch 复现了 Decoder-Only 架构的轻量级语言模型，涵盖 Pretrain、Full SFT、LoRA SFT、DPO 训练阶段。\par
\textbf{项目工作：}
\begin{enumerate}
  \item \textbf{纯 PyTorch 从零实现 Decoder-Only Transformer：}配置 d\_model=768、n\_layer=8、GQA(q=8, kv=4)、YaRN + RoPE、MoE 架构（4 experts / Top-1 路由 + 负载均衡 Loss）、Pre-Norm + RMSNorm、SwiGLU 激活、vocab=6400，参数量 198M-A64M。
  \item \textbf{完成四阶段训练链路：}Pretrain、Full SFT、垂直领域 LoRA SFT、DPO 对齐。
  \item \textbf{工程优化：}AMP 混合精度 + 梯度累积 + cosine LR 调度 + KV-Cache 推理加速。
  \item \textbf{拓展实验：}复现 MoonshotAI 团队提出的 Attention Residuals，替换传统残差连接，结果显示训练稳定，loss 与 baseline 接近，深层 block 输出幅值增长更受控制，AttnRes 在深度上学会选择性读取历史层。
\end{enumerate}

\newpage
\section{项目经历（续）}
\projectheading{小模型 Tool Use / Function Calling 能力提升研究}{2026-05 $\sim$ 至今}
\textbf{项目介绍：}工具调用是当前大语言模型落地应用的核心能力之一。然而受限于参数规模，小模型在工具调用任务上往往表现欠佳。本项目针对 Llama3.2-1B-Instruct 等小参数量模型在工具调用中常见的格式错误、工具误选、参数幻觉等痛点，基于 BFCL v4 评测框架，通过 Prompt 工程、SFT 及 RL 等手段，系统探索提升小模型工具调用能力（Single Turn、Hallucination Measurement）的有效路径。\par
\textbf{项目工作：}
\begin{enumerate}
  \item \textbf{数据准备和处理：}
    \begin{itemize}
      \item 基于 xLAM-function-calling-60k 原始语料，分类采样以控制数据集配比（40/20/30/10）。工具描述格式、输出要求格式共 8 变体（python/json），防止模型仅记忆单一格式而减弱泛化能力。
      \item 蒸馏高质量工具调用结果：以 DeepSeek-V4 为教师 A 生成推理链（<think>...</think>）与工具调用结果；以 MiniMax-m2.5 为教师 B 对推理链做蒸馏替换，对比不同教师模型在下游任务上的效益差异。
      \item 实现带重试的清洗脚本，过滤多次格式错误样本，得到 SFT 数据集（DS 2k/5k/10k 与 MM 2k/5k/10k）。
      \item 针对 irrelevance 场景利用大模型合成 300 条 no-call 负样本，覆盖无关工具、参数缺失、参数模糊／冲突、约束不支持四类子场景。
    \end{itemize}
  \item \textbf{SFT 消融实验：}
    \begin{itemize}
      \item 在 Llama 3.2-1B-Instruct 上完成七组主实验（DS 2k/5k/10k + MM 2k/5k/10k + baseline prompt-only），评测基于 BFCL v4 Single Turn、Hallucination Measurement 相关子集：Live、Non-Live。
      \item 最佳结果：MiniMax-5k 取得 Non-Live Acc 68.6\%（baseline 38.4\%，\textbf{+30.2pp}），Live Acc 51.0\%（baseline 11.4\%，\textbf{+39.6pp}）；DeepSeek-2k 取得 Overall 13.81\%（baseline 10.1\%，\textbf{+3.71pp}）。
    \end{itemize}
  \item \textbf{Badcase 分析与改进：}
    \begin{itemize}
      \item 对 MiniMax-2k 模型进行系统性 Badcase 归类，应拒用却调用的幻觉情况占比最高。原因是 SFT 训练中全部为正调用样本。
      \item MiniMax-2k 加入 300 条负样本后重训，Overall Irrelevance 从 3.48\% 提升至 53.58\%（\textbf{+50.1pp}），同时发现 trade-off（Live Acc 46.93\% $\rightarrow$ 40.04\%，Relevance Detection 100\% $\rightarrow$ 93.75\%），指明负样本比例消融实验是下一步方向。
    \end{itemize}
  \item \textbf{GRPO 训练：}
    \begin{itemize}
      \item 奖励函数设计：设计细粒度奖励函数：r\_format $\in$ [0,1]；r\_correcet = r\_name + r\_param + r\_value，归一化到 [-3,3]。
      \item 实验设置：Llama3.2-1B-Instruct 能力有限，切换为 Qwen2.5-3B-Instruct 模型；TRL 框架，vLLM 推理生成 rollout；五组实验：GRPO1500-only、SFT400 cold start + GRPO1500、SFT400、SFT1900、baseline。
      \item 主要实验结果：GRPO1500-only Non-Live 52.23\%，Live Acc 45.68\%；SFT400 + GRPO1500 Non-Live Acc 46.21\%，Live Acc 37.75\%；相比 baseline 均有所降低（61.67\%，51\%），SFT 组基本与 baseline 持平。分析原因主要来自数据集，后续可更换更高质量和多样性的数据集进行实验。
    \end{itemize}
\end{enumerate}
